\documentclass[a4paper,10pt]{article}
\usepackage[utf8]{inputenc}

% ----  Useful packages % ---- 
\usepackage{amsmath}
\usepackage{graphicx}
\usepackage{amsfonts}
\usepackage{amsthm}
\usepackage{amssymb}
\usepackage{makecell}
% ----  Useful packages % ---- 

\usepackage{wrapfig}
\usepackage{caption}
\usepackage{subcaption}
\usepackage{hyperref}
\hypersetup{
    colorlinks,
    citecolor=black,
    filecolor=black,
    linkcolor=black,
    urlcolor=black
}

\graphicspath{ {./images/} }

% ---- Set page size and margins replace ------
\usepackage[letterpaper,top=2cm,bottom=2cm,left=3cm,right=3cm,marginparwidth=1.75cm]{geometry}
% ---- Set page size and margins replace ------

% ------- NOTA ------
\theoremstyle{remark}
\newtheorem{note}{Note}[subsection]
% ------- NOTA ------

% ------- OSSERVAZIONE ------
\theoremstyle{definition}
\newtheorem{observation}{Osservazione}[subsection]
% ------- OSSERVAZIONE ------

% ------- DEFINIZIONE ------
\theoremstyle{plain}
\newtheorem{definition}{Definizione}[subsection]
% ------- DEFINIZIONE ------

% ------- ESEMPIO ------
\theoremstyle{definition}
\newtheorem{example}{Esempio}[subsection]
% ------- ESEMPIO ------

% ------- DIMOSTRAZIONE ------
\theoremstyle{definition}
\newtheorem{demostration}{Dimotrazione}[subsection]
% ------- DIMOSTRAZIONE ------

% ------- TEOREMA ------
\theoremstyle{definition}
\newtheorem{theorem}{Teorema}[subsection]
% ------- TEOREMA ------

% ------- COROLLARIO ------
\theoremstyle{plain}
\newtheorem{corollaries}{Corollario}[theorem]
% ------- COROLLARIO ------

% ------- PROPOSIZIONE ------
\theoremstyle{plain}
\newtheorem{proposition}{Proposizione}[subsection]
% ------- PROPOSIZIONE ------

% ---- Footer and header ---- 
\usepackage{fancyhdr}
\pagestyle{fancy}
\fancyhf{}
\fancyhead[LE,RO]{A.A 2023-2024}
\fancyhead[RE,LO]{Calcolo Numerico}
\fancyfoot[RE,LO]{\rightmark}
\fancyfoot[LE,RO]{\thepage}

\renewcommand{\headrulewidth}{.5pt}
\renewcommand{\footrulewidth}{.5pt}
% ---- Footer and header ---- 

% ----  Language setting ---- 
\usepackage[italian, english]{babel}
% ----  Language setting ---- 

\usepackage{listings}
\usepackage{color}

\definecolor{dkgreen}{rgb}{0,0.6,0}
\definecolor{gray}{rgb}{0.5,0.5,0.5}
\definecolor{mauve}{rgb}{0.58,0,0.82}

\lstset{frame=tb,
  language=C,
  aboveskip=3mm,
  belowskip=3mm,
  showstringspaces=false,
  columns=flexible,
  basicstyle={\small\ttfamily},
  numbers=none,
  numberstyle=\tiny\color{gray},
  keywordstyle=\color{blue},
  commentstyle=\color{dkgreen},
  stringstyle=\color{mauve},
  breaklines=true,
  breakatwhitespace=true,
  tabsize=3
}

\title{\textbf{Calcolo Numerico}}
\author{Realizzato da: Giuntoni Matteo e Filippo Ghirardini}
\date{A.A. 2023-2024}

\begin{document}
\begin{titlepage} %crea l'enviroment
	\begin{figure}[t] %inserisce le figure
		\centering\includegraphics[width=0.98\textwidth]{marchio_unipi_pant541.png}
	\end{figure}
	\vspace{20mm}
	
	\begin{Large}
		\begin{center}
			\textbf{Dipartimento di Informatica\\ Corso di Laurea Triennale in Informatica\\}
			\vspace{20mm}
			{\LARGE{Corso a Libera Scelta - 6 CFU}}\\
			\vspace{10mm}
			{\huge{\bf Computer Graphics}}\\
		\end{center}
	\end{Large}
	
	
	\vspace{36mm}
	%minipage divide la pagina in due sezioni settabili
	\begin{minipage}[t]{0.47\textwidth}
		{\large{\bf Professore:}\\ \large{Prof. }}
	\end{minipage}
	\hfill
	\begin{minipage}[t]{0.47\textwidth}\raggedleft
		{\large{\bf Autore:}\\ \large{Filippo Ghirardini}}
	\end{minipage}
	
	\vspace{25mm}
	
	\hrulefill
	
	\vspace{5mm}
	
	\centering{\large{\bf Anno Accademico 2023/2024 }}
	
\end{titlepage}
\tableofcontents
\newpage
\maketitle
\begin{center}
    \vspace{-20pt}
    \rule{11cm}{.1pt} 
\end{center}
\newpage
\section{Aritmetica di Macchina}
Per una macchina la scrittura $(x + y) + z \neq x + (y + z)$. Vediamo dunque che ci sono alcuni punti focali da considerare
per far si che una macchina funzioni correttamente:

\begin{itemize}
    \item Trovare uno standard per come memorizzare i numeri.
    \item Trovare uno standard per come manipolare i numeri.
\end{itemize}

\hspace{-15pt}Da questi due punti possiamo ricondurci ad un solo problema, come andare a rappresentare i numeri.

\subsection{Teorema di rapresentazione}
\begin{theorem}
    Dato $x \in \mathbb{R}, x \neq 0$ esistono e sono univocamente determinati.
    \begin{enumerate}
        \item un intero $p \in \mathbb{Z}$ detto esponente della rappresentazione.
        \item una successione di numeri naturiali $\{d_i\}_{i\geq 1}$ con $d_i \neq 0, 0 \leq d_i \leq B - 1$ e $d_i$
        non definitvamente uguali a \(B - 1\), dette cire della rappresentazione; tali per cui si ha
    \end{enumerate}
    \begin{equation}
        x = sign(x)B^p \sum_{i=1}^{+\infty}d_i B^{-1}.
    \end{equation}
\end{theorem}

\hspace{-15pt}Andiamo ora ad analizzare il significato di questo teorema. Esso descrive quella che viene chiamata rappresentazione
in virgola mobile, in quanto l'esponente \(p\) on è determinato in modo da avere la parte intera nulla. Le cose da considerare
in questo teorema sono:
\begin{itemize}
    \item La condizione \(d_i \neq 0\) e \(d_i\) non definitivamente uguale a \(B - 1\) sono introdotte per garantire
    l'unicità delle rappresentazioni. Ad esempio:
    \[B = 10 \text{ abbiamo } 1 = +10^1 (1 \cdot 10^{-1}) = +10^2 (0 \cdot 10^{-1} + 1 \cdot 10^{-1})\]
    Quindi due rappresentazioni diverse per lo stesso numero, però considerando le condizioni scritte sopra la seconda non risulta
    accettabile perché la prima cifra è nulla.
    \item Il caso \(x=0\) non ammette rapresentazione normalizzata. Questa casistica viene trattata dalla macchina in un modo
    particolare, per questo abbiamo la condizione \(x\neq \).
    \item Questa rapresentazione si estende anche all'insieme dei numeri complessi del tipo \(z = a + ib\), utilizzando una rapresentazione
    come coppie di numeri reali del tipo \((a,b)\).
\end{itemize}

Possiamo dedurre che visto che stiamo lavornano con registri di meoria di un calcolatore con memoria a numero finito, anche la quantità
di cifre rapresentabili saranno a numero finito esso vinene chiamato \textbf{insieme dei numeri di macchina}.\\

Dal teorema di rapresentazione in base di un numero reale può avvenire assegnando delle posizioni di meoria per il segno, 
per l'esponente e per le cifre della rappresentazione.

\begin{definition}[Insieme die numeri di macchina]
    Si definisce l'insieme dei nuermi di macchina in rappresentazione floating point con t cifre, base B e range \(-m, M\) l'insieme dei numeri reali.
    \[\mathbb{F}(B, t, m, M) = \{0\} \cup \{s \in \mathbb{R} \::\: x = sign(x) B^p \sum_{i=1}^{t} d_i B^{-1}, 0\leq d_i \leq B - 1, d_1 \neq 0, -m \leq p \leq M\}\]
\end{definition}

\hspace{-15pt}Si osserva in questa definizione che:
\begin{itemize}
    \item L'insieme \(\mathbb{F}\) ha cardinalità finita \(N = 2B^{t-1}(B-1)(M + m + 1) + 1\).
    \item L'insieme dei numeri di macchina \(\mathbb{F}(B, t, m, M)\) è simmetrico rispetto all'origine.
    \item Possiamo definire \(\Omega = B^M (B-1) \sum_{i=1}^{t} b^{-1}\) come il più grande numero macchina e \(\omega = B^{-m} B^{-1}\) come
    invece il più piccolo.
    \item Posto un \(x = B^P \sum_{i=1}^{t}d_i B^{-1}\) possiamo definire il suo successiovo numero di macchina come
    \(y = B^p(\sum_{i=1}^{t-1}d_i B^{-1} + (d_t + 1)B^{-t})\). Da qui vediamo che la distanza \(y - x = B^p - t\) porta
    i numeri ad essere non equispaziali fra di loro, quindi la distanza aumento con l'avicinarci a \(\Omega\)
\end{itemize}

\begin{example}
    Facciamo ora un esempio in cui andiamo a rappresentare il numero successivo di \(x = B^p \sum_{t}^{i=1} d_i B^{-1}\). Esso
    si può scrivere come \(y = B^p \bigg( \sum_{i=1}^{t-1} d_i B^{-1} + (d_t + 1)B^{-t} \bigg)\).\\
    Mentre si può scrivere la distanza fra questi due valori come \(y - x = B^p - t\).
\end{example}

\hspace{-15pt}E' stato fissato uno standard IEEE 754 fra gli anni 70/80, questo standard dice che, visto ci sono macchine che hanno metodi di rappresentazione
diversi bisogna fissare un standard, esso appunto dice che \(B = 2\) ed i registri sono a 32 o 64 bit.\\\\
Questa rapresentazione ha uno svantaggio che può sembrare minimo ma non lo è, lo 0 si rappresenta due volte con \(-0, +0\). Per ovviare
a questo problema si è andato ad abbandonare questa rappresentazione in esponenti ma si rapprensentato i numeri nel seguente modo:
\(p_1 2^0 + p_1 2^1 + \dots + p_11 s^10 \) che rappresentano numeri da 0 a \(2^11 - 1\) quindi 2047 numeri, mentre lo 0 si può scrivere come:
\begin{itemize}
    \item O tenendo tutti i valori a 0
    \item Oppure tendendo tutti i valori a 1
\end{itemize}

\hspace{-15pt}In entrambi i casi abbiamo un range di valori che va da \([-1022, 1024]\). A questo punto ho \(2^{P-1022}\) numeri che la 
macchina rappresneta come \(\pm 2^{P - 1022}(0.1d_1 \dots d_{52})\).\\
Impostando questo standard abbiamo \(\Omega = 2^{1024}(01 \dots 1)_2\) e \(\omega = 2^{-1022} (101)_2\).\\

\begin{observation}
    Quando \(p=0\) abbiamo i numeri che si trovano nella porzione della retta dei numeri che è compresa fra \(-\omega\) \(\omega\)
    e possiamo qui avere anche tuti 0 e quindi si introduce il caso dei numeri denormalizzati.
\end{observation}

\hspace{-15pt}Se abbiamo l'esponante uguale a tutti 1, la convensione è che tutte le cifre della mantissa sono tutti uguali a 0/1
questo numero indica il \(\pm \infty\) altrimenti sta a signifiare NaN (not a number). Questi valori ci permettono di gestire
forme indeterminate.

\subsection{Errore di rappresentazione}
Quando si va a rappresentare un numero reale non nullo \(x \in \mathbb{R}\) e con \(x \neq 0\) si può andare a commettere degli 
errori di rappresentazione detto anche \textbf{errore relativo di approssimazione}, e di definisce come, prendendo un \(\tilde{x} \in \mathbb{F}(B, t, m, M)\)
\[\epsilon_x = \frac{\tilde{x} - x}{x} = \frac{\eta x}{x}, x \neq 0\]

\hspace{-15pt}Definiamo \(|\epsilon_x| = \big | \frac{\tilde{x} - x}{x}\big | \leq \frac{B^{P-t}}{| x|} \leq \frac{B^{P - t}}{B^{P - 1}} = B^{1 - t} = u\)
la \(u\) è definita come \textbf{precisione di macchina}.\\

\hspace{-15pt}Andiamo inoltre a definire le conidizioni di underflow e overflow. Dato un \(x \in \mathbb{R}, x \neq 0\) abbiamo che:
\begin{enumerate}
    \item Se \(|x| < \omega\) o \(|x| > \Omega\) overflow. In questo caso si va ad associare il \(+\infty\).
    \item Se invece \(\omega \leq |x| \leq \Omega\) abbiamo underflow. In questo caso allora prendiamo una \(x = B^p \sum_{i = 1}^{\infty}d_i B^{-1}
    \longrightarrow B^p \sum_{i=1}^{t}d_i B^{-1} = \tilde{x}\) che è una approssimazione 
\end{enumerate}
% !TeX spellcheck = it_IT
\newpage
\section{Calcolo degli errori}
Supponiamo di avere una funzione \(f: [a,b] \to \mathbb{R}\) e \(f \neq 0\), per andare a calcolare questa funzione come già visto usiamo un algoritmo
che esprime tale valore come risutlato di una sequanza di operazioni aritmentiche. Questa rappresentazione come abbiamo già potuto
verificare con esempi produce degli errori di approssimazione. Questi errori possono essere suddivisini in 3 tipologie.

\begin{definition}[Errore inerente o inevitabile]
    Si dice errore inerente o inevitabile generato nel calcolo di \(f(x) \neq 0\) la quantità:
    \[\epsilon_{in} = \frac{f(\tilde{x}) - f(x)}{f(x)}\]
\end{definition}

\begin{definition}[Errore algoritmico]
    Si dice errore algoritmico generato nel calcolo di \(f(x) \neq 0\) la quantità:
    \[\epsilon_{alg} = \frac{g(\tilde{x}) - f(\tilde{x})}{f(\tilde{x})}\]
\end{definition}

\begin{definition}[Errore totale]
    Si dice errore algoritmico totale nel calcolo di \(f(x) \neq 0\) mediante l'algoritmo specificato da g la quantità:
    \[\epsilon_{tot} = \frac{g(\tilde{x}) - f(x)}{f(x)}\]
\end{definition}

\begin{observation}
    Vediamo che se \(|\epsilon_{in}|\) è grande il problema si definisce \textbf{problema mal condizionato}. Mentre se
    \(|\epsilon_{alg}\) è grande l'algoritmo si dice che \textbf{algoritmo è numericamente instabile}.
\end{observation}
% !TeX spellcheck = it_IT
\newpage
\section{Matrici}
%TODO recuper

\subsection{Norme vettoriali}
Sono uno strumento che ci permette di definire una distanza tra due vettori.
\begin{definition}[Norma vettoriale]
	È una funzione del tipo $f:\mathbb{R}^n \to \mathbb{R}$ che deve soddisfare tre proprietà:
	\begin{enumerate}
		\item $f(x)\geq 0 \wedge f(x)=0 \Leftrightarrow x=0$
		\item $f(\alpha x)=\lvert \alpha \rvert f(x) \quad \forall \alpha \in \mathbb{R} \quad \forall x \in \mathbb{R}^n$
		\item \textbf{Disuguaglianza triangolare}: $f(x+y)\leq f(x)+f(y) \quad \forall x,y \in \mathbb{R}^n$
	\end{enumerate}
	e la indichiamo come
	\begin{equation*}
		f(x) = \lVert x \rVert
	\end{equation*}
\end{definition}

Detto questo possiamo definire una distanza come:
\begin{equation}
	dist(x,y)=\lVert  x-y \rVert
\end{equation}
Le tre proprietà ci danno alcune informazioni sulla distanza:
\begin{enumerate}
	\item La distanza deve essere non negativa e valere $0$ solo se i due vettori coincidono 
	\item La distanza tra $x$ e $y$ deve essere uguale a quella tra $y$ e $x$
	\item $\lVert x-y \rVert = \lVert(x-a)+(a-y)\rVert \leq \lVert x-a \rVert + \lVert a-y \rVert$
\end{enumerate}

\begin{definition}[Distanza euclidea - Norma 2]
	\begin{equation}
		f(x) = \sqrt{x_1 ^ 2 + x_2 ^ 2 + \ldots + x_n ^ 2} > \lVert  x \rVert_2
	\end{equation}
\end{definition}
\begin{definition}[Norma infinito]
	\begin{equation}
		f(x)=max\lvert x_1 \rvert = \lVert x \rVert_{\infty}
	\end{equation}
\end{definition}
\begin{definition}[Norma 1]
	\begin{equation}
		f(x)=\sum_{i=1}^{n}\lvert x_i \rvert = \lVert x\rVert_1
	\end{equation}
\end{definition}

\begin{example}
	Prendiamo due vettori
	\begin{equation*}
		\begin{Bmatrix}
			1 \\
			1
		\end{Bmatrix} \quad
		\begin{Bmatrix}
			2 \\
			-2
		\end{Bmatrix}
	\end{equation*}
	e calcoliamo le varie norme:
	\begin{equation*}
		\lVert x \rVert_2 = \sqrt{2} \quad \lVert x \rVert_\infty = 1 \quad \lVert x \rVert_1 = 2
	\end{equation*}
	%TODO Finisci
\end{example}

\begin{definition}[Norma matriciale]
		È una funzione del tipo $f:\mathbb{R}^{n\times n} \to \mathbb{R}$ che deve soddisfare tre proprietà:
	\begin{enumerate}
		\item $f(A)\geq 0 \wedge f(A)=0 \Leftrightarrow A=0$
		\item $f(\alpha A)=\lvert \alpha \rvert f(x) \quad \forall \alpha \in \mathbb{R} \quad \forall A \in \mathbb{R}^{n\times n}$
		\item \textbf{Disuguaglianza triangolare}: $f(A+B)\leq f(A)+f(B) \quad \forall A,B \in \mathbb{R}^{n\times n}$
		\item $f(A\cdot B) \leq f(A) f(B)$
	\end{enumerate}
	e la indichiamo come
	\begin{equation*}
		f(A) = \lVert A \rVert
	\end{equation*}
\end{definition}
% !TeX spellcheck = it_IT
\newpage
\section{Condizionamento}
Studiare il condizionamento di un problema in forma $Ax=b$ significa chiedersi di quanto cambia la soluzione perturbando di poco $A$ e $B$.
\subsection{Metodi diretti}
Dato un sistema lineare $Ax=b$ le soluzioni possono essere trovate tramite
\begin{equation*}
	x_i = f_i(a_{11}, \ldots, a_{1n}, b1, \ldots, b_n) \quad i:1\ldots n
\end{equation*}
Si può partire studiando la forma di $A$ per cercare dei sistemi risolvibili facilmente.
\subsubsection{Matrice diagonale}
Il primo caso è quando $A$ è una matrice \textbf{diagonale}:
\begin{equation*}
	a_{ij}=0 \Longleftarrow i \neq j 
\end{equation*}
Il determinante di una matrice diagonale è $det(A)=\prod_{i=1}^{n} a_i$ e il determinante è diverso da $0$ solo se tutti gli elementi della diagonale lo sono.
Possiamo riscrivere la matrice in un sistema di equazioni lineari:
\begin{equation*}
	Ax=b \Leftrightarrow \left\{ \begin{array}{lr}
		a_1x_1 = b_1 \Leftrightarrow x1 = \frac{b1}{a1}\\
		\vdots\\
		a_nx_n = b_n \Leftrightarrow x_n=\frac{b_n}{a_n}
	\end{array} \right.
\end{equation*}
Questo sistema lo risolviamo in $O(n)$. Dato però che ridurre una matrice normale in una diagonale è un processo complesso, non conviene farlo.
\subsubsection{Matrice triangolare}
Una matrice è \textbf{triangolare inferiore} quando $a_{ij}=0$ per $j>i$, mentre è \textbf{triangolare superiore} quando $a_{ij}=0$ per $i>j$.\\
Per calcolare il determinante di una matrice triangolare superiore con Laplace facciamo:
\begin{equation*}
	A=\begin{bmatrix}
		a_{11} & 0 & \cdots & 0 \\
		\vdots & \ddots & 0 & \vdots \\
		\vdots & \ddots& \ddots & 0\\
		a_{n1} & \cdots &\cdots & a_{nn}
	\end{bmatrix}
\end{equation*}
Per risolvere $Ax=b$ con una matrice triangolare, vediamo che il sistema associato ci porta ad avere prima un equazione con tutte le incognite fino all'ultima con una sola incognita. È intuitivo che per calcolarne la soluzione è sufficiente partire dall'ultima ed eseguire una \textbf{sostituzione all'indietro}.
\begin{lstlisting}[language=Matlab]
	function[x]=backward_substitution(a,b)
		n=length(b);
		x = zeros(n,1);
		for k=n:-1:1
			s=0;
			for j=k+1:n
				s=s+a(k,j)*x(j);
			end
			x(k)=(b(k)-s)/a(k,k);
		end
\end{lstlisting}
Dato che l'operazione moltiplicativa in macchina richiede leggermente più lunga di quella della somma, è sufficiente considerare le prime per calcolare la complessità di questo programma.\\ La complessità al caso pessimo è quindi $\frac{n \cdot (n+1)}{2} = O(n^2) = \frac{n^2}{2}+O(n)$.
\begin{example}[Matrice bidiagonale superiore]
	Supponiamo di avere la seguente matrice bidiagonale superiore:
	\begin{equation*}
		A=\begin{bmatrix}
			a_{11} & a_{12} & 0 & 0\\
			0 & \ddots & \ddots & 0 \\
			0 & 0 & \ddots & a_{n-1 n} \\
			0 & 0 & 0 & a_{nn}
		\end{bmatrix}  = 
		\begin{bmatrix}
			a_1 & b_1 & 0 & 0\\
			0 & \ddots & \ddots & 0\\
			0 & 0 & \ddots & b_{n-1}\\
			0 & 0 & 0 & a_n
		\end{bmatrix}
	\end{equation*}
	In questo caso il codice si può cambiare sostituendo il secondo ciclo con l'operazione
	\begin{lstlisting}[language=Matlab]
		s=a(k,k+1)*x(k+1);
	\end{lstlisting}
	e la complessità diventa $O(n)$.
\end{example}
% !TeX spellcheck = it_IT
\section{Metodi iterativi}
Una classe di metodi che permette di risolvere sistemi lineari è quella dei metodi iterativi. L'idea è di costruire una successione di vettori partendo dalla matrice che convergano alla soluzione del sistema lineare.
\begin{equation*}
	\{X_k\}_{k\in \mathbb{N}}
\end{equation*}
\begin{observation}
	Non potendo generare infiniti termini, ad un certo punto ci sarà bisogno di fermarsi quando si pensa di essere sufficientemente vicini alla soluzione.
\end{observation}
Diciamo che
\begin{equation}
	\lim_{k\to +\infty}x_k \Leftrightarrow \lim_{k\to +\infty} \lvert\lvert x_k - x\rvert\rvert_{\infty} = 0
\end{equation}
\noindent Questo è vero perché:
\begin{equation*}
	\forall j = 1, \ldots, n \quad 0 \leq \lvert x^{(k)}_j - x_j \rvert \leq \lvert\lvert x^{(k)} - x\rvert\rvert_{\infty}
\end{equation*}

Abbiamo una matrice invertibile $A$. Supponiamo di scomporre la matrice come:
\begin{equation}
	A=M-N
\end{equation}
con l'assunzione che anche $M$ sia invertibile ($det(M)=0$). A questo punto possiamo dire che:
\begin{equation*}
	Ax=b \Leftrightarrow (M-N)x = b\Leftrightarrow Mx=Nx+b \Leftrightarrow x=M^{-1}Nx + M{-1}b \Leftrightarrow x=Px+q
\end{equation*}
Dovendo trovare la soluzione di $x=Px+q$, possiamo scegliere un vettore iniziale $x_0 \in \mathbb{R}^n$ e costruirci la successione di vettori
\begin{equation}
	\label{equation:succ}
	x_{k+1}=Px_k + q
\end{equation}
Se questa successione converge ho trovato la soluzione del sistema lineare.
\begin{observation}
	Nell'implementazione pratica non userò mai l'equazione \ref{equation:succ} ma invece
	\begin{equation}
		Mx_{k+1}=Nx_k + b
	\end{equation}
\end{observation}

\begin{example}
	Prendiamo
	\begin{equation*}
		\begin{bmatrix}
			2 & 1 \\
			1 & 2
		\end{bmatrix} \quad
		b = \begin{bmatrix}
			0 \\ 0
		\end{bmatrix} \quad
		Ax=b \Leftrightarrow x=\begin{bmatrix}
			0 \\ 0
		\end{bmatrix}
	\end{equation*}
	Partiamo definendo le due matrici per la scomposizione:
	\begin{equation*}
		M=\begin{bmatrix}
			2 & 0 \\ 0 & 2
		\end{bmatrix} \quad
		N = \begin{bmatrix}
			0 & -1 \\ -1 & 0
		\end{bmatrix}
	\end{equation*}
	e calcolando la matrice $P$
	\begin{equation*}
		P = M^{-1}N = \begin{bmatrix}
			\frac{1}{2} & 0 \\ 0 & \frac{1}{2}
		\end{bmatrix}\begin{bmatrix}
		0 & -1 \\ -1 & 0
		\end{bmatrix} = \begin{bmatrix}
		0 & -\frac{1}{2} \\ -\frac{1}{2} & 0
		\end{bmatrix}
	\end{equation*}
	e il vettore $q$
	%TODO
	Iniziamo il metodo iterativo:
	\begin{equation*}
		\begin{split}
			& x^{k+1} = Px^{(k)} = \begin{bmatrix}
				0 & -\frac{1}{2} \\ -\frac{1}{2} & 0
			\end{bmatrix} x^{(k)} \\
			& x^{(1)} =  \begin{bmatrix}
				0 & -\frac{1}{2} \\ -\frac{1}{2} & 0
			\end{bmatrix} \begin{bmatrix}
			x \\ y
			\end{bmatrix} = \begin{bmatrix}
			-\frac{1}{2} y \\ -\frac{1}{2} x
			\end{bmatrix} \\
			& x^{(2)} =  \begin{bmatrix}
				0 & -\frac{1}{2} \\ -\frac{1}{2} & 0
			\end{bmatrix} \begin{bmatrix}
				-\frac{1}{2}y \\ -\frac{1}{2}x
			\end{bmatrix} = \begin{bmatrix}
				\frac{1}{4} x \\ \frac{1}{4} y
			\end{bmatrix}
		\end{split}
	\end{equation*}
	e notiamo che la successione tende a 
	\begin{equation*}
		\begin{bmatrix}
			0 \\ 0
		\end{bmatrix}
	\end{equation*}
	Se prendiamo invece
	\begin{equation*}
		M=\begin{bmatrix}
			0 & 1 \\ 1 & 0
		\end{bmatrix} \quad
		N = \begin{bmatrix}
			-2 & 0 \\ 0 & -2
		\end{bmatrix}
	\end{equation*} con
	\begin{equation*}
		P = M^{-1}N = \begin{bmatrix}
			0 & 1 \\ 1 &0
		\end{bmatrix}\begin{bmatrix}
			-2 & 0 \\ 0 & -2
		\end{bmatrix} = \begin{bmatrix}
			0 &-2 & -2 & 0
		\end{bmatrix}
	\end{equation*}
	abbiamo che la successione \textbf{diverge}.
\end{example}
\begin{definition}
	Un metodo iterativo è convergente se
	\begin{equation}
		\forall x^{(0)} \in \mathbb{R}^n \quad x_k \to x
	\end{equation}
	Ovvero se per ogni vettore di partenza scelto, il metodo converge.
\end{definition}

\begin{theorem}
	Dato $x^{(k+1)} = Px^{(k)}+q$, se
	\begin{equation}
		\exists \lvert\lvert\cdot\rvert\rvert \text{ tc } \lvert\lvert P \rvert\rvert <1
	\end{equation}
	allora il metodo è convergente.
\end{theorem}

\begin{example}
	Data una matrice
	\begin{equation*}
		A = \begin{bmatrix}
			3 & -1 & \\ 
			-1 & \ddots & \ddots \\
			 & \ddots & \ddots & -1\\
			 & & -1 & 3
		\end{bmatrix}
	\end{equation*}
	definiamo le matrici per la scomposizione
	\begin{equation*}
		M = \begin{bmatrix}
			3 \\
			& \ddots \\
			& & 3
		\end{bmatrix} \quad
		N = \begin{bmatrix}
			0 & 1 \\
			1 & \ddots & \ddots \\
			& \ddots & \ddots & 1 \\
			& & 1 &0
		\end{bmatrix}
	\end{equation*}
	Per capire se il metodo è convergente dobbiamo calcolare la norma infinito di $P$:
	\begin{equation*}
		M^{-1} N = \begin{bmatrix}
			0 & \frac{1}{3} \\
			\frac{1}{3} & \ddots & \ddots \\
			& \ddots & \ddots & \frac{1}{3}\\
			& & \frac{1}{3} & 0
		\end{bmatrix} \quad
		\lvert\lvert P \rvert\rvert_{\infty}= \frac{1}{3} + \frac{1}{3} = \frac{2}{3} < 1
	\end{equation*}
	Il problema di questo metodo iterativo è che ha complessità $O(n)$ per ogni iterazione e non è quindi competitivo con l'eliminazione Gaussiana.
\end{example}

\begin{example}
	Data una matrice
	\begin{equation*}
		A = \begin{bmatrix}
			T & I & \\ 
			-I & \ddots & \ddots \\
			& \ddots & \ddots & I\\
			& & -I & T
		\end{bmatrix} \quad
		T = \begin{bmatrix}
			5 & -1 & \\ 
			-1 & \ddots & \ddots \\
			& \ddots & \ddots & -1\\
			& & -1 & 5
		\end{bmatrix}
	\end{equation*}
	In questo caso i metodi iterativi sono vantaggiosi poiché anche se la matrice è predominante diagonale (e quindi permette Gauss senza problemi), a causa dell'effetto del fill-in lo rende sconveniente.
\end{example}

\begin{example}
	Data una matrice
	\begin{equation*}
		A = \begin{bmatrix}
			n & -1 & \ldots & -1\\ 
			-1 & \ddots & \ddots & \vdots \\
			\vdots & \ddots & \ddots & -1\\
			-1 & \ldots & -1 & n
		\end{bmatrix}
	\end{equation*}
	calcoliamo $P$
	\begin{equation*}
		N = \begin{bmatrix}
			0 & 1 & \ldots & 1\\ 
			1 & \ddots & \ddots & \vdots \\
			\vdots & \ddots & \ddots & 1\\
			1 & \ldots & 1 & 0
		\end{bmatrix} \quad
		P = \begin{bmatrix}
			0 & \frac{1}{n} & \ldots & \frac{1}{n}\\ 
			\frac{1}{n} & \ddots & \ddots & \vdots \\
			\vdots & \ddots & \ddots & \frac{1}{n}\\
			\frac{1}{n} & \ldots & \frac{1}{n} & 0
		\end{bmatrix}
	\end{equation*}
	La sua norma vale
	\begin{equation*}
		\lvert\lvert P \rvert\rvert_{\infty} = \frac{n-1}{n}
	\end{equation*} e quindi converge.
\end{example}

\begin{note}
	Se la norma vale $1$ non posso dire nulla sulla convergenza.
\end{note}

\begin{theorem}
	\begin{equation}
		x^{(k+1)} = Px^{(k)}+q \text{ è convergente } \Leftrightarrow \phi(P)<1
	\end{equation}
\end{theorem}

\begin{example}
		Data una matrice
	\begin{equation*}
		A = \begin{bmatrix}
			1 & & & \alpha\\ 
			-1 & \ddots &  & \\
			& \ddots & \ddots &\\
			& & -1 & 1
		\end{bmatrix}
	\end{equation*}
	troviamo le matrici per la scomposizione
	\begin{equation*}
		M=I \quad
		N = \begin{bmatrix}
			0 & &  & -\alpha\\ 
			1& \ddots & &\\
			 & \ddots & \ddots &\\
			&& 1 & 0
		\end{bmatrix}
	\end{equation*}
	Per quali $\alpha$ il metodo è convergente? \\
	Troviamo la fattorizzazione LU della matrice per il calcolo del determinante e otteniamo così il polinomio caratteristico:
	\begin{equation*}
		x^n + \alpha
	\end{equation*}
	Dobbiamo poi trovare il modulo degli autovalori per poterli confrontare:
	\begin{equation*}
		\begin{split}
			x^n=-\alpha \\
			\lvert\lambda\rvert^n = \lvert -\alpha \rvert \Rightarrow \lvert \lambda \rvert = \sqrt[n]{\lvert \alpha \rvert} \\
			\sqrt[n]{\lvert \alpha \rvert} < 1 \Leftrightarrow \lvert \alpha \rvert < 1
		\end{split}
	\end{equation*}
\end{example}
\end{document}
